% Generated draft: requires mathematical and source review.
\documentclass[11pt]{article}
\usepackage[margin=25mm]{geometry}
\usepackage{fontspec}
\setmainfont{Latin Modern Roman}
\usepackage{amsmath,amssymb,mathtools,unicode-math}
\setmathfont{Latin Modern Math}
\usepackage{enumitem,needspace,xurl,hyperref}
\setlength{\parindent}{0pt}
\setlength{\parskip}{5pt}
\setlength{\emergencystretch}{3em}
\newcommand{\N}{\mathbb{N}}
\newcommand{\Z}{\mathbb{Z}}
\newcommand{\Q}{\mathbb{Q}}
\newcommand{\R}{\mathbb{R}}
\newcommand{\C}{\mathbb{C}}
\newcommand{\F}{\mathbb{F}}
\newcommand{\A}{\mathbb{A}}
\newcommand{\PP}{\mathbb{P}}
\newcommand{\E}{\mathbb{E}}
\newcommand{\eps}{\varepsilon}
\newcommand{\dd}{\,\mathrm{d}}
\begin{document}
\section*{UnsolvedMath ID: 30000259}
\textbf{Source problem code:} catalogue entry
\subsection{Definitions and mathematical statement}
The record asks four independent extensions of the braid Conley index. In the intended setting, fix a skeleton braid $y$ and a relative class $[x\,\mathrm{rel}\,y]$ of free strands, with the properness and boundedness/admissibility conditions required by the theory. A discretization with $d$ sampling steps gives a space $\mathcal C_d([x\,\mathrm{rel}\,y])$ of admissible discrete representatives. The exact strand category, boundary/periodicity convention and sampling rules are not supplied.

For orientation, a Conley index of an isolated invariant set of a specified flow can be represented by the pointed homotopy type $N/L$, where $(N,L)$ is an index pair (an isolating neighborhood with its appropriate exit subset). The braid Conley index applies this construction in the source's braid-class setting; a generic homotopy quotient without isolation and the specified flow is not a substitute.
\begin{enumerate}
\item Which braid-group data are sufficient to compute the braid Conley index? This asks for an effective relation to the relevant braid information; a word, conjugacy class, relative skeleton data and further admissibility information are not automatically the same input.
\item What is the smallest admissible $d$ for which a relative braid class has just one component in the discretized space? With path-component convention this asks for nonemptiness and $|\pi_0(\mathcal C_d)|=1$. The original mentions a ``stated minimal discretization'', but its proposed numerical bound or formula is absent, so none is guessed; existence and behavior under further subdivision are not asserted.
\item When does the index yield dynamical information beyond the braid class alone? The intended comparison may concern forcing invariant sets, equilibria or periodic behavior, but the source gives no precise strengthening criterion. This should not be rephrased as two identical braid classes having different topological invariants without stating the dynamics and comparison convention.
\item Does the theory extend to general fourth-order conservative equations? The equation class, ODE/PDE setting, boundary conditions and meaning of conservative (for example energy-preserving or Hamiltonian) are not defined. An extension must specify replacement dynamics, isolation/admissibility and an invariant with the intended properties; the ordinary second-order/parabolic braid construction cannot simply be presumed unchanged.
\end{enumerate}
The supplied literature triage checked 21 August 2026 reports substantial later braid-Conley-index work but no joint resolution of all four questions. It specifically retains the discretized-component conjecture and general conservative fourth-order extension as unresolved in that triage. It supplies no complete later-paper references or independent current status for each of the other parts.
\subsection{Short English statement}
The entry asks how braid data compute the index, when discretization yields one relative-class component, what extra dynamical conclusions the index gives, and whether the theory extends to fourth-order conservative equations. All four are preserved, including the missing discretization bound and equation-class hypotheses.
\subsection{Sources}
\begin{itemize}
\item Local UnsolvedMath registry, permanent ID 30000259, exact source code catalogue entry; \emph{Extensions of the Braid Conley Index}. Proposed year 2005, status open. The corrected\_verified statement restores all four questions independently; its claimed source inspection is metadata, not an inspection performed here.
\item Supplied report citation: \emph{Topological and Variational Methods for Differential Equations (2006)}, Oberwolfach Reports, DOI 10.4171/owr/2005/29, pp. 1601--1678, \url{https://doi.org/10.4171/owr/2005/29}. Title year differs from DOI/proposed year.
\item Supplied literature triage checked 2026-08-21, citing only ``original problem report'', \url{https://doi.org/10.4171/owr/2005/29}. It reports subsequent work but retains the component and conservative-fourth-order questions unresolved. The dated OWR-derived curation requests independent verification of time-sensitive status and cited claims; no separate later reference is supplied.
\end{itemize}
\end{document}
